%% File: chapters/chapter6.tex
%% Author: Y.U.P. (yashparitkar)
%%
%% Description: Driving the UART wrapper from the command line.

\documentclass[../manual.tex]{subfiles}

\begin{document}

\chapter{The Command Line Interface}
\label{ch:cli}

\todo{Every command with its usage and a worked scenario. The transcripts come from real runs against hardware, not from plausible-looking output written by hand.}

\section{Telling the tool about a core}
\label{sec:cli_devices}

\todo{\texttt{--add-device} connects, reads the UID back and writes the device store; after that a core is addressed by \texttt{--device}. The point to make loudly: \texttt{--device} selects by \emph{UID}, not by the order devices were added, and a core built without a G\_UID reports 0. \texttt{--device-dir} has to agree between the CLI and the GUI.}

\section{The verbs and the order they run in}
\label{sec:cli_verbs}

\todo{The flags are verbs, and they are applied in a fixed order whatever order they are given in. The ordering table, then a whole capture as one line, and the same capture split across invocations: nothing about the session is kept on the host.}

\section{Configuring the trigger}
\label{sec:cli_trigger}

\todo{Condition and mask as one number covering the whole probe word, in any base, with bit 0 the first signal in the portmap. These do \emph{not} use the portmap names. Type and reduction, omitted fields keeping what the core already had, and a value with bits above the probe width being refused rather than trimmed.}

\section{Worked scenarios}
\label{sec:cli_scenarios}

\todo{Three or four end-to-end runs with real pasted output: trigger on a condition, force a trigger that never came, disarm and retry, and a capture split across two invocations.}

\section{Debugging the link}
\label{sec:cli_debug}

\todo{\texttt{--debug} logs every UART transaction. A real log extract and how to read it: the byte counts matter as much as the errors, because a reply that is short or long for its request means the two ends disagree about which request is being served.}

\section{Exit statuses}
\label{sec:cli_exit}

\todo{The exit codes, a sentence each on what actually causes them, and a pointer to the troubleshooting chapter.}
\end{document}